\chapter{Validation du prototype}

\section{Stratégie de test}

La validation combine des tests unitaires et des essais sur la plateforme
complète. Les tests unitaires vérifient les règles du simulateur et le format
des notifications. Les essais d'intégration publient de vrais messages dans
Mosquitto, passent par Node-RED et contrôlent les points écrits dans InfluxDB.
Enfin, les requêtes Flux de chaque panneau sont exécutées par l'API de Grafana.

\section{Tests unitaires}

Douze tests automatisés sont présents :

\begin{itemize}
  \item huit tests du simulateur : identité des messages, date effective de
        péremption, consommation en attente, capacités des fluides, codes et
        paramètres 94xx, débit de marquage et absence de marquage sans demande ;
  \item quatre tests du notifier : format d'une alerte active, regroupement,
        résolution et protection contre un template Grafana non résolu.
\end{itemize}

La commande \texttt{make verify} exécute ces tests, vérifie la syntaxe Compose,
le JSON des flux Node-RED et les fichiers des tableaux de bord.

\section{Essais d'intégration}

\begin{table}[H]
\centering
\caption{Résultats des essais de bout en bout}
\begin{tabularx}{\textwidth}{p{0.31\textwidth}Yp{0.16\textwidth}}
\toprule
Essai & Observation & Résultat \\
\midrule
Publication simulée &
Télémétrie, événements et marquages présents dans le bucket. &
Conforme \\
Message QoS 1 publié deux fois &
Le nombre de points reste inchangé après la retransmission exacte. &
Conforme \\
Payload sans \texttt{message\_id} &
Message reçu sur \texttt{sgx/system/dead-letter} avec le motif. &
Conforme \\
Redémarrage Node-RED &
Reconnexion MQTT puis reprise des écritures. &
Conforme \\
Authentification Node-RED &
\texttt{/flows} renvoie HTTP 401 sans session ; les quatre onglets sont
accessibles après connexion. &
Conforme \\
Source Grafana &
Connexion InfluxDB déclarée saine et bucket accessible. &
Conforme \\
Tableaux de bord &
25 requêtes de panneaux exécutées avec succès. &
Conforme \\
Alertes &
Sept règles chargées depuis les fichiers de provisioning. &
Conforme \\
\bottomrule
\end{tabularx}
\end{table}

\begin{resultbox}
Au terme des essais, les six conteneurs étaient déclarés sains. Les trois
mesures InfluxDB contenaient des points issus du simulateur, et la répétition
d'un message conservait le même résultat. Le prototype valide donc la chaîne
simulateur--MQTT--Node-RED--InfluxDB--Grafana.
\end{resultbox}

\section{Cas de la disponibilité}

Le simulateur alterne les périodes de commande et les périodes vides. Les
échantillons sont classés dans quatre catégories : production, arrêt réel,
attente avec jet actif et absence de commande. Le tableau de bord ne compte
que la deuxième catégorie comme arrêt.

Cet essai a révélé une limite connue. La requête quotidienne compte des
échantillons et suppose une publication périodique. Ce calcul est correct pour
le simulateur. Si Litmus Edge publie seulement lors d'un changement de valeur,
les états longs et courts auraient le même poids. La validation sur site devra
donc vérifier le mode de publication.

\section{Cas des consommables}

Le lot de solvant initial est ouvert depuis 85 jours, avec une durée
configurée de 90 jours après ouverture. Sa date imprimée est plus éloignée,
mais la requête calcule cinq jours restants. Ce scénario vérifie que la plateforme
retient la première échéance.

Les remplissages produisent un saut de niveau. Lorsque la hausse dépasse cinq
points, Node-RED réinitialise le débit estimé. Le délai avant épuisement reste
alors vide jusqu'à la réception d'une nouvelle baisse mesurable, au lieu
d'afficher une prévision négative ou héritée de l'ancienne cartouche.

Les tests de conformité 94xx contrôlent aussi la capacité de 800~mL des
cartouches externes, la capacité de 1,075~L des réservoirs internes et les codes
du jet. Un état \texttt{running} avec demande produit ainsi un code jet 7 et un
état d'impression 1, tandis que la phase de démarrage produit le code jet 1.

\section{Limites de la validation}

Le prototype ne prouve pas encore la compatibilité avec les tags réels de
l'imprimante. Il ne mesure pas non plus la latence du réseau industriel ni la
qualité d'un payload Sparkplug B fourni par le site. Les défauts simulés sont
plausibles, mais leurs codes ne remplacent pas le manuel du modèle installé.

Ces limites ne remettent pas en cause les composants en aval. Elles définissent
le travail de raccordement : capturer un payload réel, établir la table de
mapping, vérifier les unités et rejouer les tests d'acceptation.
